%! AP Statistics — Unit 8, Lesson 8.6 — Homework (Answer Key)
\documentclass[10pt]{article}
\usepackage{apstats-article}
\usepackage{apstats-key}   % pulls in -boxes; do NOT also load -boxes
\newcommand{\TallMath}[1]{$\displaystyle #1\rule[-1.4em]{0pt}{3.2em}$}

\begin{document}

\pageheader{Unit 8, Lesson 8.6}{Homework --- Answer Key}
\namedateperiod

\vspace{0.08in}

\textit{Show all work. Use all four SPDC steps for Problem 1.}

\begin{enumerate}[itemsep=10pt]

%% Problem 1 %%
\item \textbf{Sleep and Grade Level.}

\renewcommand{\arraystretch}{1.4}
\begin{center}
\begin{tabular}{@{} l c c c @{}}
\hline
           & \textbf{$<$7 hrs} & \textbf{7$+$ hrs} & \textbf{Total} \\
\hline
9th grade  & 30 & 30 & 60  \\
10th grade & 38 & 22 & 60  \\
11th grade & 42 & 18 & 60  \\
\hline
\textbf{Total} & 110 & 70 & 180 \\
\hline
\end{tabular}
\end{center}

\begin{enumerate}[label=\alph*., itemsep=6pt]
  \item \textbf{State.}
        Test type: \ans{Homogeneity --- three separate groups (grade levels) each randomly sampled; one variable (sleep amount).}

        $H_0$: \ansline{The distribution of sleep amount is the same across all three grade levels.}
        $H_a$: \ansline{The distribution of sleep amount differs for at least one grade level.}

  \item \textbf{Plan.} Method: \ans{Chi-square test for homogeneity.}
        Random: \ansline{Random selection stated in the problem.}

        $E_{\text{9th, }<7} = \dfrac{60 \times 110}{180} = $ \ans{36.67}
        \quad $E_{\text{9th, }7+} = \dfrac{60 \times 70}{180} = $ \ans{23.33}

        \ans{All three grade levels have the same expected counts (36.67 and 23.33), all $\ge 5$. Condition met.}

  \item \textbf{Do.} $\text{df} = ($ \ans{3} $-1)($ \ans{2} $-1) = $ \ans{2}

        \renewcommand{\arraystretch}{1.6}
        \begin{center}
        \begin{tabular}{@{} l c c @{}}
        \hline
                   & \textbf{$<$7 hrs} & \textbf{7$+$ hrs} \\
        \hline
        9th grade  & $\dfrac{(30-{\color{keyred}\mathbf{36.67}})^2}{{\color{keyred}\mathbf{36.67}}} \approx 1.21$
                   & $\dfrac{(30-{\color{keyred}\mathbf{23.33}})^2}{{\color{keyred}\mathbf{23.33}}} \approx 1.91$ \\
        10th grade & \ans{$\dfrac{(38-36.67)^2}{36.67} \approx 0.05$} & \ans{$\dfrac{(22-23.33)^2}{23.33} \approx 0.08$} \\
        11th grade & \ans{$\dfrac{(42-36.67)^2}{36.67} \approx 0.78$} & \ans{$\dfrac{(18-23.33)^2}{23.33} \approx 1.22$} \\
        \hline
        \end{tabular}
        \end{center}

        $\chi^2 = 1.21 + 1.91 + 0.05 + 0.08 + 0.78 + 1.22 \approx $ \ans{5.25}

        \begin{teachernote}
        Exact: $(30-36.\overline{6})^2/36.\overline{6} + \cdots$. With exact fractions $\chi^2 = 5.25$; rounding intermediate steps may give 5.25--5.33. Using the table: df = 2, $\chi^2 = 5.25$ falls between 4.605 ($p = 0.10$) and 5.991 ($p = 0.05$), so the p-value is between 0.05 and 0.10. Accept $p \approx 0.07$ from technology.
        \end{teachernote}

        p-value: \ans{Between 0.05 and 0.10 (from table, df = 2: $\chi^2 = 5.25$ lies between 4.605 and 5.991); technology gives $p \approx 0.072$.}

  \item \textbf{Conclude.} At $\alpha = 0.05$:
        \ansline{Because $p \approx 0.07 > \alpha = 0.05$, we fail to reject $H_0$. We do not have convincing evidence that the distribution of sleep amount differs across grade levels.}
        \ansline{}
\end{enumerate}

\vspace{0.1in}

%% Problem 2 %%
\item \textbf{Degrees of Freedom.}

\begin{enumerate}[label=\alph*.]
  \item $\text{df} = (2-1)(4-1) = $ \ans{3}
  \item Critical value at $\alpha = 0.05$, df $= 3$: \ans{7.815} (from table)
  \item \ansline{$\chi^2 = 9.21 > 7.815$, so yes, the student can reject $H_0$ at $\alpha = 0.05$. The p-value is less than 0.05 (in fact, between 0.025 and 0.05).}
\end{enumerate}

\vspace{0.1in}

%% Problem 3 %%
\item \textbf{Conceptual Question.}

\ansline{A larger $\chi^2$ value means the observed counts deviate more from the expected counts we would predict if $H_0$ were true. These larger deviations are less likely to occur by random chance alone. In terms of the chi-square distribution, a larger test statistic falls farther in the right tail, resulting in a smaller p-value. A smaller p-value means the data provide stronger evidence against $H_0$, making it more likely that the variables are truly associated (or the distributions differ).}
\ansline{}
\ansline{}
\ansline{}

\end{enumerate}

\end{document}
